\documentclass{article}
\usepackage{amsmath, amssymb, amsthm}

\setlength{\parindent}{0pt}

\newtheorem{claim}{Claim}

\begin{document}

The army needs to identify soldiers with a disease called ``klep". There is a way to test blood to determine 
whether it came from someone with klep. The straightforward approach is to test each soldier individually. 
This requires $n$ tests, where $n$ is the number of soldiers. 

\medskip

A better approach is the following: group the soldiers into groups of $k$. Blend the blood samples of each 
group and apply the test once to each blended sample. If the group blend doesn't have klep, we are done with 
that group after one test. If the group blend fails the test, then someone in the group has klep, and we 
individually test all the soldiers in the group.

\medskip

Assume each soldier has klep with probability, $p$, independently of all the other soldiers.

\begin{claim}
    Assuming for simplicity that $n$ is divisible by $k$, the expected number of tests as a function of $n$, 
    $p$, and $k$ is
    \[
    \frac{n}{k} + n(1 - (1-p)^k)
    \]
\end{claim}

\begin{proof}
    For each $i \in \{1, 2, \dots, \frac{n}{k}\}$, define the random indicator variable
    \[
    X_i =
    \begin{cases}
        1 & \text{if group } i \text{ has at least one soldier with klep} \\
        0 & \text{otherwise}
    \end{cases}
    \]
    The probability that group $i$ has no soldiers with klep is $(1-p)^k$. Thus,
    \[
    \mathrm{E}[X_i] = 1\left(1 - (1-p)^k\right) + 0\left((1-p)^k\right) = 1 - (1-p)^k
    \]
    Let $T$ denote the number of tests for all groups. Each group requires one initial test, and if $X_i = 1$, 
    it requires an additional $k$ tests. Therefore,
    \[
    T = \sum_{i=1}^{n/k}(1 + kX_i)
    \]
    By linearity of expectation,
    \[
    \begin{aligned}
        \mathrm{E}[T] &= \sum_{i=1}^{n/k}\left(1 + k\mathrm{E}[X_i]\right) \\[4pt]
        &= \frac{n}{k}\left(1 + k\left(1 - (1-p)^k\right)\right) \\[4pt]
        &= \frac{n}{k} + n(1 - (1-p)^k)
    \end{aligned}
    \]

    \medskip

    Hence the expected number of tests is 
    \[
    \frac{n}{k} + n(1 - (1-p)^k)
    \]
\end{proof}

\begin{claim}
    The expected number of tests is approximately minimised by choosing 
    \[
    k \approx \frac{1}{\sqrt{p}}
    \]
    and the resulting expected number of tests is 
    \[
    \mathrm{E}[T] \approx 2n\sqrt{p}
    \]
\end{claim}

\begin{proof}
    Since $n$ is a positive scalar independent of $k$, minimising $\mathrm{E}[T] = n \cdot g(k)$ over $k$ is 
    equivalent to minimising
    \[
    g(k) = \frac{1}{k} + 1 - (1-p)^k.
    \]
    Treating $k$ as a continuous positive real variable and differentiating with respect to $k$,
    \[
    g'(k) = -\frac{1}{k^2} - (1-p)^k \ln(1-p)
    \]
    Setting $g'(k) = 0$ and rearranging gives
    \[
    \frac{1}{k^2} = -(1-p)^k \ln(1-p)
    \]
    Since $0 < p < 1$, we have $0 < 1-p < 1$, so $\ln(1-p) < 0$. Hence the right-hand side of the equation is
    positive, consistent with the left-hand side.

    \medskip

    The above equation is difficult to solve explicitly for $k$, since $k$ appears in both a power-law term 
    and an exponential term. We therefore approximate by assuming $p \ll 1$. Using the first-order Taylor 
    expansions,
    \[
    \ln(1-p) \approx -p \quad \text{and} \quad (1-p)^k = e^{k\ln(1-p)} \approx e^{-kp}
    \]
    Hence the equation becomes 
    \[
    \frac{1}{k^2} \approx pe^{-kp}
    \]
    Further assuming that $kp \ll 1$ and applying $e^{-kp} \approx 1 - kp$, we have 
    \[
    \frac{1}{k^2} \approx p(1 - kp) = p - kp^2
    \]
    Assuming for now that $kp^2 \ll p$, we discard it, leaving 
    \[
    \frac{1}{k^2} \approx p \quad \Longrightarrow \quad k \approx \frac{1}{\sqrt{p}}.
    \]
    Substituting the resulting estimate $k \approx p^{-1/2}$ gives
    \[
    kp \approx \sqrt{p} \ll 1 \quad \text{and} \quad kp^2 \approx p^{3/2} = o(p)
    \]
    confirming that both approximations are self-consistent.

    \medskip

    Now it remains to compute the resulting expected number of tests. Substituting $k = \frac{1}{\sqrt{p}}$ 
    into $g(k)$, 
    \[
    g(k) \approx \sqrt{p} + 1 - (1 - \sqrt{p}) = 2\sqrt{p}
    \]
    where the last step again uses $e^{-kp} \approx 1 - kp$, with $kp \approx \sqrt{p} \ll 1$. Thus
    \[
    \mathrm{E}[T] = n \cdot g(k) \approx 2n\sqrt{p}
    \]

    \medskip

    Hence the expected number of tests is approximately minimised by choosing 
    \[
    k \approx \frac{1}{\sqrt{p}}
    \]
    and the resulting expected number of tests is 
    \[
    \mathrm{E}[T] \approx 2n\sqrt{p}
    \]
\end{proof}

\begin{claim}
    The fraction of the work that the grouping method saves over the straightforward approach is approximately 
    $80\%$ in a million-strong army where $1\%$ have klep.
\end{claim}

\begin{proof}
    Using the approximate formula $\mathrm{E}[T] \approx 2n\sqrt{p}$, the fraction of work saved is
    \[
    1 - \frac{\mathrm{E}[T]}{n} \approx 1 - 2\sqrt{p} = 1 - 2\sqrt{0.01} = 1 - 0.2 = 0.80
    \]

    \medskip

    Verifying with the exact formula, we first calculate $k = \frac{1}{\sqrt{p}}$ for $p = 0.01$, giving 
    $k = 10$. Substituting this into the exact formula from Claim~1 gives
    \[
    \mathrm{E}[T] = \frac{n}{k} + n\left(1-(1-p)^k\right) = \frac{10^6}{10} + 10^6\left(1-0.99^{10}\right) 
    \approx 195{,}618
    \]
    so the fraction of work saved is
    \[
    1 - \frac{195{,}618}{10^6} \approx 0.804
    \]
    Since $0.804 \approx 0.80$, the exact formula agrees closely with the approximate formula.

    \medskip

    Hence the grouping method saves $80\%$ of the work over the straight-forward approach.
\end{proof}

\end{document}
